% Figure: the power triangle. Average power P (real axis), reactive
% power Q (imaginary axis), apparent power S as the hypotenuse, and the
% power-factor angle phi between P and S. A second dashed triangle shows
% the effect of power-factor correction.
\begin{figure}[htbp]
\centering
\begin{tikzpicture}[scale=1.0, line cap=round, >=latex]
  \draw[->, engblack, thin] (-0.4, 0) -- (5.2, 0) node[right, font=\small] {$P$ (W)};
  \draw[->, engblack, thin] (0, -0.4) -- (0, 3.4) node[above, font=\small] {$Q$ (var)};
  % Uncorrected triangle: P = 4, Q = 3, S = 5, phi = atan(3/4) = 36.87 deg.
  \draw[engblue, very thick, ->] (0, 0) -- (4, 3);
  \filldraw[engblue] (4, 3) circle (0.05);
  \node[engblue, font=\footnotesize, anchor=south west] at (4, 3) {$S$};
  \draw[engblue, thick] (0, 0) -- (4, 0);
  \draw[engblue, thick] (4, 0) -- (4, 3);
  \node[engblue, font=\footnotesize, anchor=north] at (2, 0) {$P$};
  \node[engblue, font=\footnotesize, anchor=west] at (4, 1.5) {$Q$};
  % power-factor angle
  \draw[engblue] (0.8, 0) arc (0:36.87:0.8);
  \node[engblue, font=\footnotesize] at (1.05, 0.27) {$\varphi$};
  % Corrected triangle: same P = 4, smaller Q' = 1.3, S' shorter, phi' smaller.
  \draw[engred, very thick, ->] (0, 0) -- (4, 1.3);
  \filldraw[engred] (4, 1.3) circle (0.05);
  \node[engred, font=\footnotesize, anchor=south west] at (4, 1.25) {$S'$};
  \draw[engred, dashed, thin] (4, 3) -- (4, 1.3);
  \node[engred, font=\footnotesize, anchor=west] at (4.05, 2.15) {$\Delta Q$ (capacitor)};
  \draw[engred] (1.2, 0) arc (0:18.0:1.2);
  \node[engred, font=\footnotesize] at (1.45, 0.2) {$\varphi'$};
\end{tikzpicture}
\caption[The power triangle]{The power triangle places average power
$P$ on the real axis, reactive power $Q$ on the imaginary axis, and
apparent power $S = P + iQ$ as the hypotenuse. The angle $\varphi$
between $P$ and $S$ is the power-factor angle; the power factor is
$\cos\varphi = P/|S|$. Power-factor correction adds a capacitor that
supplies leading reactive power $\Delta Q$, shrinking the load's
reactive power from $Q$ to $Q'$ (red) at the same real power $P$. The
corrected apparent power $S' < S$ reduces the current the supply must
deliver, which is why the supply transformer and conductors can be
sized smaller after correction.}
\label{fig:vol02:ch03:power-triangle}
\end{figure}
